%\section{Related Work}

%\paragraph{LLM sycophancy} Sharma et al.~\cite{sharma2025understandingsycophancylanguagemodels} definesycophancy as the tendency for instruction-tuned models to agree with user views even when incorrect, with Wei et al.~\cite{wei2023} demonstrating this effect across model sizes. Subsequent work expands beyond factual agreement to social adaptation~\cite{cheng2025} and preference mirroring. These works study how LLMs are influenced by user preferences and beliefs, while we study how LLMs are influenced by an action being framed as being theirs.

%\paragraph{LLM self-bias} Models favor their own outputs over identical content from others~\cite{spiliopoulou2025playfavoritesstatisticalmethod}, or prefer other models’ outputs to those of humans~\cite{laurito2025}. These studies focus on \textit{different answers}—comparing alternative completions and showing preference for the option closer in style or content to the model itself. Self-evaluation bias specifically has been measured by Koo et al.~\cite{koo2023} and Panickssery et al.~\cite{panickssery2024llmevaluatorsrecognizefavor}, who document preference for self-generated content. In contrast, we study how models evaluate the \textit{same answer} differently depending on context. We find evidence of \textit{self-attribution bias}, where models downplay risks of actions they think they have generated.

\section{Related Work}

\paragraph{LLM Self-Preference in Evaluation.}
Recent work establishes that LLMs systematically favor their own outputs. \citet{panickssery2024llmevaluatorsrecognizefavor} showed GPT-4 recognizes its own generations and demonstrated a link between self-recognition and self-preference. \citet{wataoka2024selfpreferencebias} traced this effect to perplexity-based familiarity: models assign higher scores to lower-perplexity text, which it's own outputs necessarily are, resulting in emergent self-preference. \citet{spiliopoulou2025playfavoritesstatisticalmethod} found this extends to family-level bias, with models preferring outputs from architecturally similar systems.  

\textbf{Conceptually, self-preference concerns which actions or outputs a model favors among alternatives. In contrast, self-attribution bias concerns how a model’s judgment about the \emph{same action} changes solely due to attribution context, such as whether the action is framed as the model’s own.}

Prior work characterizes self-preference as a stylistic or familiarity-based bias operating over static artifacts. Self-attribution bias, by contrast, arises when models evaluate actions they have actually produced and that are explicitly framed as their own, exhibiting commitment effects—such as larger bias when evaluation follows generation—that cannot be explained by distributional similarity alone.





% These findings characterize self-preference as a stylistic bias operating through implicit familiarity detection on static artifacts.Self-attribution bias, by contrast, arises from on-policy generation: when models evaluate actions they have actually produced, explicitly framed as their own, they exhibit commitment effects like larger bias when evaluation follows generation.

\paragraph{Biases in LLM-as-Judge Systems.}
A growing literature examines LLMs as evaluators, with implications for their use as safety monitors. \citet{zheng2023judgingllmasajudge} documented position bias, verbosity bias, and self-enhancement effects. \citet{wang2023largelanguagemodelsare} showed rankings can be manipulated simply by reordering responses. \citet{chen2025dollmevaluatorsprefer} used verifiable benchmarks to show that self-preference causes evaluation failures disproportionately when models are wrong. \citet{tsui2025selfcorrectionbench} documented a ``self-correction blind spot'' where models correct errors from external sources while failing to correct identical errors in their own outputs--direct evidence that attribution modulates evaluation. We extend these findings to agentic safety contexts, showing self-attribution bias is largest on incorrect code and harmful computer use actions, precisely when reliable self-monitoring is most needed.

\paragraph{Implicit Collusion and Coordination without Intent.}
Our findings are also related to concerns about collusion and coordination failures in agentic and oversight systems.
Prior work has studied explicit or strategic collusion, in which multiple agents intentionally coordinate to evade monitoring or pursue shared objectives \citep{hubinger2024scheming,carlsmith2023collusion}. Related research in scalable oversight emphasizes that evaluator independence is critical, as coupling generators and judges can lead to correlated failures even without adversarial intent \citep{irving2018ai,kenton2021alignment,bowman2022failure}. Separately, work on self-consistency and confidence calibration shows that agreement among a model’s own outputs can amplify both correct and incorrect beliefs \citep{wang2022self}.

In contrast to these settings, the behavior we document does not involve intentional coordination, communication, or shared objectives. 
Instead, self-attribution bias induces a form of \emph{implicit collusion}: when the same model (or tightly coupled instances of a model) is reused for action and evaluation, risk signals are systematically suppressed and errors are reinforced rather than corrected.
This produces outcomes analogous to collusion: overly permissive approvals, reduced error detection, and inflated confidence. 


%\paragraph{Biases in LLM-as-a-judge evaluations}
%LLM evaluators exhibit systematic biases including position~\cite{wang2024a}, verbosity, and familiarity effects~\cite{stureborg2024largelanguagemodelsinconsistent}. Self-evaluation bias specifically has been measured by Koo et al.~\cite{koo2023} and Panickssery et al.~\cite{panickssery2024llmevaluatorsrecognizefavor}, who document preference for self-generated content. 
%Critically, existing methods conflate self-preference with quality differences and lack frameworks to isolate bias from confounds~\cite{xu2025}. We address this through calibrated metrics that account for position/label effects, revealing that bias intensifies after action commitment -- not merely preferring but rationalizing risky choices post-hoc.

%While prior work establishes evaluation biases exist, no study examines how \textit{commitment to actions} affects subsequent evaluations. This distinction is safety-critical: a model that prefers its outputs is concerning; a model that downplays risks after acting on them is dangerous. We demonstrate this commitment bias represents a distinct failure mode requiring targeted mitigation.
